% ============ 02. Contexto y guía de lectura ============
\section{Contexto y guía de lectura}

\subsection{Qué es este documento}

Este documento desarrolla el \textbf{criterio 4 de la rúbrica}:
\emph{Validación de la calidad de pruebas, tratamientos y lógica estadística}. El repositorio auditado es el
Observatorio de investigadores reconocidos del Ministerio de Ciencia,
Tecnología e Innovación de Colombia (MinCiencias), que consolida y analiza
\textbf{6 convocatorias (2013--2021)} del padrón oficial de investigadores
y su producción científica. Esta serie cubre los \textbf{6 criterios} de la
rúbrica en documentos separados, más el documento maestro que los integra.

\textbf{En resumen:} Este criterio evalúa el rigor estadístico: cómo se limpiaron los datos (atípicos, faltantes) y si los métodos usados (matrices de transición, índices de concentración, comparaciones de género y diversidad, redes, agrupamientos) son correctos. Conclusión: la estadística descriptiva es sólida y transparente; la principal debilidad es la falta de inferencia formal (intervalos de confianza, pruebas de supuestos, modelos ajustados).




\begin{tcolorbox}[nota]
\textbf{Cómo leer este documento:} cada PDF de esta serie desarrolla \emph{un}
criterio de la rúbrica. Los demás criterios se tratan a fondo en sus propios
documentos y en el \textbf{documento maestro} (\texttt{maestro\_auditoria.pdf}).
Si este es tu primer acercamiento, lee primero la portada y este contexto; al
final encontrarás las conclusiones del criterio.
\end{tcolorbox}

\subsection{Glosario (solo los términos que usa este documento)}

Esta tabla incluye únicamente los términos técnicos que aparecen en este
documento, explicados en lenguaje sencillo:

\begingroup
\small
\setlength{\tabcolsep}{3pt}
\begin{longtable}{>{\raggedright\arraybackslash}p{4.3cm}>{\raggedright\arraybackslash}p{9.8cm}}
\caption{Glosario de términos presentes en este documento}
\label{tab:glosario}\\
\toprule
\textbf{Término} & \textbf{Qué significa} \\
\midrule
\endfirsthead
\toprule
\textbf{Término} & \textbf{Qué significa} \\
\midrule
\endhead
    HHI (Herfindahl-Hirschman) & Índice que mide si algo está concentrado o repartido. Va de 0 (todo repartido) a 1 (todo en un solo lugar). En el informe se usa para ver si los investigadores se concentran en pocos departamentos.
 \\
    Representación relativa & Cuántas veces aparece un grupo (p. ej. afrocolombianos) en el padrón frente a lo que correspondería según su peso en la población colombiana (DANE). Un valor de 0.3 indica que aparece 3 veces menos de lo esperado.
 \\
    k-prototypes & Algoritmo de agrupamiento (clustering) que agrupa personas con perfiles similares cuando los datos mezclan números (edad) y categorías (área, género).
 \\
    Pipeline & Cadena de procesos que lleva los datos desde la fuente hasta el análisis (ingesta, limpieza, transformación, modelado). También llamada 'cadena de procesamiento' o 'flujo de datos'.
 \\
    Markovianidad & Propiedad de un proceso en el que la probabilidad de pasar al siguiente estado solo depende del estado actual (y no de toda la historia). Se usa para validar si las matrices de transición de categorías son adecuadas.
 \\
    Modelos de supervivencia (Kaplan-Meier / Cox) & Técnicas estadísticas para analizar 'tiempo hasta que ocurre un evento' (p. ej. cuánto tarda un investigador en salir del sistema), incluyendo datos incompletos.
 \\
    Intervalo de confianza (IC) y bootstrap & El intervalo de confianza indica el rango plausible de una estimación (p. ej. 0.15-0.19). El bootstrap es un método que lo calcula remuestreando los datos muchas veces, sin fórmulas complejas.
 \\
\bottomrule
\end{longtable}
\endgroup
